\usepackage{amsmath, amssymb, amsthm}       % 美國數學學會
\usepackage[CheckSingle, CJKmath]{xeCJK}    % 中文
\usepackage{fontspec}                       % 字型配置
\usepackage{geometry}                       % 版面配置
\usepackage[nocheck]{fancyhdr}              % 頁首頁尾
\usepackage{color}                          % 顏色
\usepackage[x11names]{xcolor}               % 更多顏色
\usepackage{datetime2}                      % 日期、時間
\usepackage{listings}                       % 顯示 code 用的
\usepackage{titlesec}                       % 標題格式控制
\usepackage{tikz}                           % 畫圖
\usepackage{enumerate}                      % 列舉環境
\usepackage{enumitem}                       % 自訂列舉環境
\usepackage{ulem}                           % 文字底線、刪除線
\usepackage{graphicx}                       % 插入圖片
\usepackage{multicol}                       % 多欄排版
\usepackage{hyperref}                       % 超連結
\usepackage{subcaption}                     % 子圖
\usepackage{cancel}                         % 取消線
\usepackage{multirow}                       % 表格合併列
\usepackage{float}                          % 浮動物件位置強制
\usepackage{notomath}                       % Noto 數學字體

%%%%%%%%%%%%%%%%%%%%%%%%%%%%%%

\setmainfont{Noto Serif}  % 主要字體

\setmonofont{Menlo}  % 等寬字體

\setCJKmainfont{Noto Serif TC}[  % 中文主要字體
    ItalicFont     = *-Regular, % 中文沒有斜體
    BoldItalicFont = *-Bold,    % 中文沒有斜體
]

\setCJKmonofont{Noto Sans TC}[  % 中文等寬字體
    ItalicFont     = *-Regular, % 中文沒有斜體
    BoldItalicFont = *-Bold,    % 中文沒有斜體
]

\newCJKfontfamily{\kaifont}{Kaiti TC}

%%%%%%%%%%%%%%%%%%%%%%%%%%%%%%

% 彩色模式
\newcommand{\keywordcolor}{\color{Blue1}}
\newcommand{\identifiercolor}{\color{black}}
\newcommand{\commentcolor}{\color{Red4}}
\newcommand{\stringcolor}{\color{Green4}}

% % 黑白模式
% \renewcommand{\keywordcolor}{\color{black}}
% \renewcommand{\identifiercolor}{\color{black}}
% \renewcommand{\commentcolor}{\color{black!70}}
% \renewcommand{\stringcolor}{\color{black!55}}

\makeatletter
\lst@CCPutMacro\lst@ProcessOther {"2D}{\lst@ttfamily{-{}}{-{}}}
\@empty\z@\@empty
\makeatother

\lstset{
    % language=C,                             % Code 的語言
    numbers=none,                           % 行號的位置
    numberstyle=\footnotesize,              % 行號的字型和大小
    stepnumber=1,                           % 顯示行號的間隔
    numbersep=5pt,                          % 行號跟 code 的距離
    showspaces=false,                       % 顯示空白 (使用特別的底線記號)
    showtabs=false,                         % 顯示 Tab
    showstringspaces=false,                 % 顯示字串的空白
    tabsize=2,                              % Tab 的寬度
    frame=leftline,                         % adds a frame around the code
    captionpos=b,                           % sets the caption-position to bottom
    breaklines=true,                        % sets automatic line breaking
    breakatwhitespace=false,                % sets if automatic breaks should only happen at whitespace
    escapeinside={\%*}{*)},                 % if you want to add a comment within your code
    morekeywords={*},                       % 自訂的 Keywords
    basicstyle=\footnotesize\ttfamily,                   % 基本的字型和大小
    backgroundcolor=\color{white},          % 背景顏色
    keywordstyle=\bfseries\keywordcolor,    % Keywords 的字型
    identifierstyle=\identifiercolor,       % 標示符的字型
    commentstyle=\itshape\commentcolor,     % 註解的字型
    stringstyle=\itshape\stringcolor,       % 字串的字型
    % otherkeywords={mat, vec, vector},       % 其他的 Keywords
}

%%%%%%%%%%%%%%%%%%%%%%%%%%%%%%

% ===== 頁面設定 =====
\geometry{a4paper}                      % A4 紙
% \geometry{includehead}                  % 內文區塊包含頁首，沒有頁尾、旁注
\geometry{headsep=6mm}                  % 頁首與內文距離稍微拉開
\geometry{vmargin=1in, hmargin=1in}     % 上下左右邊界 1 吋

% ===== 行距與段落 =====
\renewcommand{\baselinestretch}{1.2}    % 正文行距 1.2 倍
\setlength{\parskip}{1.5ex}             % 段落間距略增，讓段落清楚分隔
% \setlength{\parindent}{0pt}             % 段落首行縮排，可依喜好調整

% ===== Chapter -> Section 風格 =====
\titleformat{\chapter}
{\normalfont\Large\bfseries}   % 原 section 大小
{\thechapter}
{1em}
{\thispagestyle{fancy}}

\titlespacing{\chapter}{0pt}{-24pt}{8pt}

\renewcommand{\chaptermark}[1]{%
    \markboth{\thechapter.\ #1}{}%
}

% ===== Section -> Subsection 風格 =====
\titleformat{\section}
{\normalfont\large\bfseries}   % 原 subsection 大小
{\thesection}
{1em}
{}

\titlespacing{\section}{0pt}{8pt}{6pt}

% ===== Subsection -> Subsubsection 風格 =====
\titleformat{\subsection}
{\normalfont\normalsize\bfseries}  % 原 subsubsection 大小
{\thesubsection}
{1em}
{}

\titlespacing{\subsection}{0pt}{6pt}{4pt}

% ===== 表格列高 =====
% \renewcommand{\arraystretch}{1.5}       % 表格列高 1.5 倍
\setlength{\tabcolsep}{1.25ex}          % 表格欄寬


% \titlespacing{\section}{0pt}{12pt}{6pt}
% \titlespacing{\subsection}{0pt}{6pt}{4pt}
% \titlespacing{\subsubsection}{0pt}{4pt}{2pt}

\XeTeXlinebreaklocale "zh"          % 中文自動換行
\XeTeXlinebreakskip = 0pt plus 1pt  % 設定段落之間的距離


%%%%%%%%%%%%%%%%%%%%%%%%%%%%%%

\hypersetup{
    colorlinks=true,         % 啟用顏色連結
    linkcolor=red,           % 內部連結顏色
    filecolor=green,         % 檔案連結顏色
    urlcolor=blue,           % 網址顏色
    pdfborder={0 0 1}        % 設定邊框樣式（0 0 1 表示只有底部有邊框）
}

%%%%%%%%%%%%%%%%%%%%%%%%%%%%%%

\newtheorem{theorem}{Theorem}
\newtheorem{corollary}[theorem]{Corollary}
\newtheorem{lemma}[theorem]{Lemma}
\newtheorem{definition}[theorem]{Definition}
\newtheorem{conjecture}[theorem]{Conjecture}

% 數域
\newcommand{\reals}{\mathbb{R}}
\newcommand{\complex}{\mathbb{C}}
\newcommand{\integers}{\mathbb{Z}}
\newcommand{\naturals}{\mathbb{N}}

% 科學記號
\newcommand{\sci}[2]{#1 \times 10^{#2}}

% 常用符號
\newcommand{\norm}[1]{\left\| #1 \right\|}
\newcommand{\abs}[1]{\left| #1 \right|}
\newcommand{\inner}[2]{\left\langle #1, #2 \right\rangle}

% 向量、矩陣
\newcommand{\mat}[1]{\boldsymbol{#1}}
\newcommand{\identity}{\mat{I}}
\newcommand{\zero}{\mat{0}}
\newcommand{\jacobian}{\mat{J}}

\newcommand{\eps}{\varepsilon}

\newcommand{\bigO}[1]{O\left( #1 \right)}

%%%%%%%%%%%%%%%%%%%%%%%%%%%%%%

% 允許跨頁的多行公式
\allowdisplaybreaks

% 調整數學模式下的行距
\setlength{\abovedisplayskip}{8pt}
\setlength{\belowdisplayskip}{8pt}
\setlength{\abovedisplayshortskip}{4pt}
\setlength{\belowdisplayshortskip}{4pt}


%%%%%%%%%%%%%%%%%%%%%%%%%%%%%%

\makeatletter
\renewcommand{\maketitle}{
    \thispagestyle{empty} % 第一頁不要頁眉
    \begin{center}
        \kaifont

        \vspace*{\fill}

        {\fontsize{40}{48} \selectfont \@title \par}

        \vspace*{\fill}

        {\large
            姓名：黃俊源 \\
            學號：01257128 \\
            日期：\@date \par
        }

        \newpage
    \end{center}
}
\makeatother

% quote 環境改用楷體
\makeatletter
\renewenvironment{quote}
{\list{}{\rightmargin\leftmargin}%
    \item\relax\kaifont}  % 加上 \kaifont
{\endlist}
\makeatother

%%%%%%%%%%%%%%%%%%%%%%%%%%%%%%
